
\subsubsection{3.1 Experimental Design}

A within-subject design was employed where the same 304 runtime error cases were repaired under all five granularity conditions. This design eliminates case-difficulty confounds.

\subsubsection{3.2 Granularity Level Definitions}

\begin{table}[htbp]
\centering
\resizebox{\columnwidth}{!}{%
\begin{tabular}{lll}
\toprule
Level & Definition & Example \\
\midrule
G0 & Pass/fail only & ''The code failed the test.'' \\
G1 & Error type & ''The code failed with: IndexError'' \\
G2 & Error + message & ''IndexError: list index out of range'' \\
G3 & Error + line & ''IndexError at line 7: list index out of range'' \\
G4 & Full stack trace & Complete Python traceback \\
\bottomrule
\end{tabular}
}
\end{table}

\subsubsection{3.3 Model and Configuration}

\begin{table}[htbp]
\centering
\resizebox{\columnwidth}{!}{%
\begin{tabular}{ll}
\toprule
Parameter & Value \\
\midrule
Model & CodeLlama-7B-Instruct \\
Source & meta-llama/CodeLlama-7b-Instruct-hf \\
Temperature & 0 (deterministic) \\
Max tokens & 512 \\
Execution timeout & 10 seconds \\
\bottomrule
\end{tabular}
}
\end{table}

Temperature 0 was used to ensure reproducibility. The same Self-Debug style prompt template was used across all conditions, with only the error feedback section varying.

\subsubsection{3.4 Statistical Analysis}

\begin{itemize}
\item \textbf{Primary analysis}: One-way ANOVA across five granularity levels
\item \textbf{Post-hoc comparisons}: Tukey's Honestly Significant Difference (HSD)
\item \textbf{Paired comparisons}: McNemar's test for G0 vs G3 and G3 vs G4
\item \textbf{Confidence intervals}: Wilson score intervals for proportions
\item \textbf{Effect size}: Eta-squared for ANOVA
\end{itemize}
